\subsection{Resultados importados}\label{sec:importados}
En esta lista \(m\) es el tamaño residual y \(n\) el orden inicial. En un residual
ordinario alcanzable, \(p\) denota el punto estacionario interior de Bethe en su soporte,
con masas uno en cada fila y columna. Escribimos \(\ell_f=\sum_{e\in f}p_e\),
\(h_f=1-\ell_f\) en ambas familias oblicuas. Las cotas comunes son
\begin{equation}\label{eq:clase-importada}
 \max_e p_e\le C/m,\qquad \min_f h_f\ge\eta>0,\qquad
 \lambda_2(\mathcal L_p)\ge\gamma>0 .
\end{equation}
El grafo bipartito de filas y columnas es conexo. Cambiar el signo de las coordenadas de
columna transforma su laplaciano \(\mathcal L_p\) en la matriz de incidencia axial ponderada
por \(p\), con núcleo \((\mathbf1_R,-\mathbf1_C)\). La brecha de~\cite{pc} y la
de~\cite{tc} son, por tanto, la misma. En las fórmulas que usan
\(\tau_e=p_e(1-p_e)/(1-4p_e)\), se tiene \(p_e\le\tau_e\le2p_e\) para \(m\) grande.
La comparación de formas conserva una brecha uniforme. El predictor emplea el cofactor
axial, el factor local y el término de normalización \(\tfrac12\log(2\pi)\).

\paragraph{Inserción de un hijo y suma por fila~\cite[teo.~7.1 y cor.~7.2]{pc}.}
Si \(m\ge\delta_0n\) y se cumple~\eqref{eq:clase-importada}, existen \(C_*,m_0\)
que sólo dependen de \(\delta_0,C,\eta,\gamma\). Para \(m\ge m_0\), todo hijo legal
tiene un único punto estacionario en la rama \(p<1/4\) y
\[
 L(B+e)-L(B)=\log[p_e(1+\delta_e)]+r_e,\qquad
 |\delta_e|\le C_*/m,\quad |r_e|\le C_*/m^2,\quad \sum_{e\in i}p_e\delta_e=0 .
\]
La corrección \(\delta_e\) es la derivada proyectada del logaritmo del prefactor dividida
por \(p_e\). El corolario da \(|\log\rho(B,i)|\le C_*/m^2\) para toda fila \(i\),
sumando todos sus hijos. Estas dos estimaciones intervienen en las identidades de historia
e inserción de los apéndices~\ref{app:signed} y~\ref{app:interior}.

\paragraph{Clase mesoscópica~\cite[\S8.1 y teo.~8.1]{pc}.}
En la misma geometría, supóngase~\eqref{eq:clase-importada} y
\(\max_f\ell_f\le\epsilon\le\epsilon_0\). Los umbrales \(\epsilon_0,m_0\) y la constante
\(C_*\) dependen de \(C,\eta,\gamma\), independientemente de \(n/m\) y del número de líneas.
Para \(m\ge m_0\) todos los hijos legales tienen su punto estacionario pequeño y
\(|\log\rho(B,i)|\le C_*/m^2\). La prueba conserva la representación anterior por hijo,
de donde \(\nu(e\mid i)=p_e(1+O(m^{-1}))\), también en la transición de primera salida.
Se usa en los apéndices~\ref{app:signed}, \ref{app:arrival} y~\ref{app:interior}.

\paragraph{Expansión del prefactor~\cite[\S\S7.5--7.8 y 8.4]{pc}.}
En soporte fijo de la clase macroscópica, con el calibre anterior,
\(s_{2,a}=\sum_{i\in a}p_i^2\) y \(\kappa_i=a_i^{\mathsf T}S_0^+a_i-4\), se usa
\[
 \ell+\tfrac12\log\det S_{0,\rm gf}+3\sum_i p_i^2
   =-14\sum_i p_i^3-\tfrac32\sum_i p_i^2\kappa_i
       +\tfrac94\sum_a s_{2,a}^2+O(m^{-2}).
\]
El resto tiene constantes de clase, con dependencia en \(\delta_0\) en el §7.
Las cotas por masa del §8.4 las hacen independientes de \(n/m\) en la clase mesoscópica.
Se usa en el apéndice~\ref{app:interior} para identificar \(mV\).

\paragraph{Parada y permanencia macroscópica~\cite[def.~2.2, \S4 y cor.~4.12]{tc}.}
Fijados \(M=\lceil\delta n\rceil\) y un suelo \(K_0\le K\le M\), hasta \(M\) se exige
que la discrepancia de intervalos, fibras y rectángulos sea menor que el radio macroscópico.
En \(M\) se añaden las tolerancias de entrada del §5.3 y la longitud máxima
\(D_{\max}(M)<\alpha_{\rm hist}M\) de los cruces vacíos. Después se exige la clase
interior de~\eqref{eq:clase-importada} con \(\ell_f\le\epsilon\), hasta \(K\) inclusive.
La primera infracción detiene la historia. Para una fracción fija \(\delta>0\), radio
suficientemente pequeño, partición suficientemente fina y \(n\) grande, la probabilidad
bajo \(\nu\) de alcanzar el radio hasta \(M\) es \(\le C_{\delta,r}e^{-c_{\delta,r}n}\).
En su complemento el máximo de Bethe es único y satisface~\eqref{eq:clase-importada}
uniformemente en las líneas. El §4.3, (4.3), da
\(mp_e/W(e)=1+O_\delta(\zeta+1/m)\) en sitios disponibles y la misma cota de error
absoluto para las masas diagonales respecto de \(1-u_f/U_f\), donde \(\zeta\) suma
discrepancias de intervalos y fibras, paso de partición y error de discretización.
A un número fijo de líneas, el Lema 4.10 asigna riesgo en un paso
\(m^{-1}(\#\text{ejes}+\sum_f(1-u_f/U_f))+O_\delta((\zeta+1/m)/m)\) bajo la ley que
elige según \(p\) en la fila. El cociente por hijo anterior conserva esa precisión bajo \(\nu\).
Usamos el corolario en el apéndice~\ref{app:signed} y el seguimiento a corte auxiliar
fijo en los apéndices~\ref{app:interior} y~\ref{app:transfer}.

\paragraph{Entrada, prescripciones y bandas~\cite[\S5, lemas 5.2, 5.7 y prop.~5.6, 5.8]{tc}.}
Se fijan los márgenes interiores y exteriores con
\(2\epsilon=\epsilon_{\rm out}\le\epsilon_0(C_{\rm out},\eta_{\rm out},\gamma_{\rm out})\),
y las tolerancias según el §5.4. El \(\delta\) de entrada debe ser suficientemente pequeño
para esas tolerancias. El Lema 5.7 aporta la entrada de los observables y una excepción
geométrica \(\le C_\lambda e^{-\lambda M}\) para \(D_{\max}(M)\).
La Proposición 5.6 reconstruye la clase interior en \(M\) y en los tableros que satisfacen
las tolerancias. Todo hijo de padre regular, incluso el de salida, hereda la clase exterior.
En ella valen los márgenes más débiles y la masa diagonal \(\le2\epsilon\) en el tramo
mesoscópico. Así los padres están en las clases de los dos teoremas de inserción,
con \(m>M\) y \(m\le M\), respectivamente. La ley \(\nu\) da a la supervivencia de
\(s\) filas y \(t\) columnas prescritas hasta \(m\), con padres regulares hasta el paso de
llegada, probabilidad \(\le(m/n)^{s+t}\exp(C(st+s+t)/m)\), uniformemente en \(s,t\).
Para todo \(K_0\le K\le M\) y \(u=2^qK\le M\), \(q\ge0\) entero, la entrada regular
seguida de primera salida en \([u,2u)\) cuesta \(\le C e^{-c_*u}\),
con \(c_*>2\beta\), \(\beta=2\epsilon_{\rm out}/(1-\epsilon_{\rm out})\).
Las constantes de bandas son comunes a los suelos. Estos son los insumos históricos de
los apéndices~\ref{app:signed} y~\ref{app:arrival}.

\paragraph{Coste estático~\cite[lemas 6.5--6.6, prop.~6.7 y cor.~6.8]{tc}.}
Sea \(\widehat W_e=p_e(1-p_e)\), y sea \(\mu_{\widehat W}\) la ley de permutaciones
del soporte con pesos \(\prod_{e\in\sigma}\widehat W_e\). El Lema 6.5, para un punto
estacionario interior con masas axiales uno y holguras positivas, da la identidad
\[
 Q(B)e^{-H_B(p)}=\frac{\operatorname{per}\widehat W}{\prod_e(1-p_e)}
                  \mathbb E_{\mu_{\widehat W}}\mathcal A,\qquad
 \mathcal A(\sigma)=
 \frac{\prod_f h_f}{\prod_e(1-p_e)^2}\,\mathbf1_{\{\sigma\ {\rm legal}\}}
 \prod_{e\in\sigma}\frac{(1-p_e)^2}{h_{D(e)}h_{A(e)}} .
\]
\(\Gamma_4,\Gamma_2\) designan los normalizadores canónicos de reinas y de torres con
los pesos \(p\) del mismo tablero. Si \(Q(B)>0\), \(p_e\le C/m\) y
\(\|p-\mathbf1\mathbf1^{\mathsf T}/m\|_{\rm op}\le\rho_*<1\), se obtiene
\(\log Q(B)-L(B)=\log\mathbb E_{\mu_{\widehat W}}\mathcal A
-\log(\Gamma_4/\Gamma_2)+o(1)\), uniformemente para \(C,\rho_*\) fijos.
La matriz axial tiene espectro \(1\pm\sigma_i(p)\), por lo que
\eqref{eq:clase-importada} implica esta última hipótesis.

Con los márgenes de~\eqref{eq:clase-importada}, el Lema 6.6 da
\((S(t)^+)_{uv}=\mathbf1_{\{u=v\}}+O_{C,\eta,\gamma}(m^{-1})\), uniforme en los
índices axiales y en \(0\le t\le1\). Aquí \(S(t)=(1-t)S_{R,2}+tS_{R,4}\), donde
\(S_{R,d}\) es la matriz de incidencia axial ponderada por \(p(1-p)/(1-dp)\).
Con \(\max_f\ell_f\le\epsilon_{\rm out}<1\), la Proposición 6.7 usa esa inversa y da
\(\log(\Gamma_4/\Gamma_2)=O_{C,\eta,\gamma}(\epsilon_{\rm out}+m^{-1})\).
El Corolario 6.8 concluye en una primera salida mesoscópica
\((\log Q-L)^+\le C_0+\beta m\), tomando cero si \(Q=0\), uniformemente en la clase exterior.
Su \emph{prueba} sirve en toda llegada de esa clase. En efecto,
\(\sum_f\log h_f-2\sum_e\log(1-p_e)\le O_C(1)\), pues
\(\sum_f\ell_f=2m=2\sum_e p_e\) y \(\sum_ep_e^2\le C\). Cada una de las \(m\)
casillas elegidas aporta a lo sumo \(2\epsilon_{\rm out}/(1-\epsilon_{\rm out})\)
al logaritmo del último producto. Los lemas 6.5--6.6 y la Proposición 6.7 absorben
el resto en \(C_0\), independientemente de \(n/m\). El apéndice~\ref{app:signed} aplica el
corolario en las bandas y el apéndice~\ref{app:arrival} aplica el argumento estático
en la llegada regular con una línea larga.

